% =============================================================================
% Sequential Pattern Analysis for Detecting Sophisticated Cheating
% in Online Gomoku: A Structural Approach to Random Engine Mixing
% Version 1 — Initial Draft
% =============================================================================
\documentclass[11pt,a4paper]{article}

% --- Packages ---
\usepackage[utf8]{inputenc}
\usepackage[T1]{fontenc}
\usepackage{lmodern}
\usepackage{amsmath,amssymb,amsthm}
\usepackage{graphicx}
\usepackage{booktabs}
\usepackage{hyperref}
\usepackage[margin=2.5cm]{geometry}
\usepackage{algorithm}
\usepackage{algpseudocode}
\usepackage{enumitem}
\usepackage{caption}
\usepackage{subcaption}
\usepackage{xcolor}
\usepackage{float}
\usepackage{natbib}

% --- Theorem environments ---
\newtheorem{definition}{Definition}
\newtheorem{proposition}{Proposition}
\newtheorem{theorem}{Theorem}
\newtheorem{lemma}{Lemma}
\newtheorem{remark}{Remark}
\newtheorem{corollary}{Corollary}

% --- Custom commands ---
\newcommand{\E}{\mathbb{E}}
\newcommand{\Prob}{\mathbb{P}}
\newcommand{\R}{\mathbb{R}}
\newcommand{\CAC}{\mathrm{CAC}}
\newcommand{\CUSUM}{\mathrm{CUSUM}}
\newcommand{\ACF}{\mathrm{ACF}}
\DeclareMathOperator*{\argmax}{arg\,max}
\DeclareMathOperator*{\argmin}{arg\,min}

% =============================================================================
\title{
    Sequential Pattern Analysis for Detecting Sophisticated Cheating\\
    in Online Gomoku: A Structural Approach to Random Engine Mixing
}

\author{
    Nguyen Cong Minh \\
    \texttt{nguyenminhfptu@gmail.com}
}

\date{\today}

% =============================================================================
\begin{document}

\maketitle

% =============================================================================
\begin{abstract}
Move-independent Bayesian frameworks for detecting engine-assisted play in
online board games have demonstrated strong performance against full-game
engine users but exhibit a systematic blind spot against \emph{partial}
or \emph{random-mixing} cheaters---players who interleave engine consultations
with genuine human moves. Such adversarial strategies produce move sequences
that, analysed move-by-move, are individually plausible yet collectively
anomalous. This paper presents a complementary detection framework grounded
in \emph{sequential pattern analysis}, targeting the structural signatures
that random mixing imprints on a move sequence as a whole.

We formalise three distinct cheating archetypes---Type~A
(complexity-targeted), Type~B (stochastic random), and Type~C
(noise-camouflaged)---and derive, for each archetype, the precise
statistical deformation it induces on the joint distribution of the
sequence $(\Delta_1, C_1), \ldots, (\Delta_n, C_n)$, where $\Delta_i$
denotes the winrate loss at move $i$ and $C_i$ denotes positional complexity.
Four detection tests are developed:
(i) the \emph{Wald--Wolfowitz Runs Test} applied to the binary-encoded
optimality sequence, grounded in combinatorics;
(ii) the \emph{Lag-1 Autocorrelation Test} (ACF), which detects the
zig-zag pattern induced by alternating engine and human moves;
(iii) the \emph{Cumulative Sum Control Chart} (CUSUM), which identifies
the \emph{change-point} at which engine use begins within a game;
(iv) the \emph{Complexity-Accuracy Correlation} (CAC), which captures
the inversion of the natural human relationship between positional
difficulty and move quality.
We further propose a \emph{Shannon Entropy test} that detects
the bimodal spike at $\Delta \approx 0$ characteristic of mixed strategies.

Asymptotic distributions, minimum sample sizes, and detection power
analyses are derived for each test under both hypotheses. An
\emph{ensemble detector} combining all five signals via a weighted
voting rule is presented, with formal analysis of the trade-off between
false-positive rate and detection power. The framework is designed
as a second-tier companion to the Bayesian per-move analysis of
\citet{nguyen2025bayesian}, forming a two-layer detection architecture.

\medskip
\noindent
\textbf{Keywords:} Anti-cheating, sequential analysis, Gomoku, random engine
mixing, Wald--Wolfowitz runs test, autocorrelation, CUSUM, change-point
detection, Shannon entropy, complexity-accuracy correlation, online gaming
integrity.
\end{abstract}

% =============================================================================
\section{Introduction}
\label{sec:introduction}

The problem of detecting engine-assisted play in online board games has been
studied primarily through a \emph{move-quality} lens: given the distribution
of winrate losses $\Delta_i$ across a game, does the aggregate evidence favour
a human or an engine user? The Bayesian framework of \citet{nguyen2025bayesian}
exemplifies this approach, modelling each $\Delta_i$ as an independent draw
from a parametric family and accumulating evidence via sequential posterior
updates.

This independence assumption is both a strength and a weakness. It is a
strength because it makes the framework analytically tractable and enables
real-time per-move updates. It is a weakness because it renders the system
\emph{blind to structure in the sequence itself}. A sophisticated cheater who
alternates between engine consultation and genuine play produces a sequence
$\Delta_1, \Delta_2, \ldots, \Delta_n$ that is not i.i.d.---it contains
systematic inter-move dependencies that a move-independent test cannot exploit.

\begin{figure}[h]
\centering
\fbox{\parbox{0.85\textwidth}{
\small
\textbf{Motivating observation.}
Consider a 22-move sequence where 40\% of moves are played by an engine
($\Delta \approx 0$) and 60\% are genuine human moves.
The \emph{aggregate} $\lambda_{\mathrm{MLE}}$ may fall within a normal
human range, causing a move-independent test to classify the game as
``Human.'' Yet the sequence contains a pronounced zig-zag pattern
(Lag-1 autocorrelation $\approx -0.1$), an anomalous spike of 64\%
moves with $\Delta \leq 1\%$, and a negative complexity-accuracy
correlation---all invisible to per-move analysis.
}}
\end{figure}

\subsection{Contributions}

This paper makes the following contributions:

\begin{enumerate}[label=(\roman*)]
\item A formal taxonomy of three sophisticated cheating archetypes
  (Types A, B, C) with precise statistical characterisations
  (Section~\ref{sec:archetypes}).

\item Four theoretically-grounded detection tests: Runs Test
  (Section~\ref{sec:runs}), ACF Test (Section~\ref{sec:acf}),
  CUSUM (Section~\ref{sec:cusum}), and CAC (Section~\ref{sec:cac}),
  each with asymptotic distributions and minimum sample size requirements.

\item A Shannon Entropy test exploiting the bimodal fingerprint of
  mixed strategies (Section~\ref{sec:entropy}).

\item An ensemble detector combining all five tests with formal analysis
  of detection power as a function of engine-use rate $p$ and sample
  size $n$ (Section~\ref{sec:ensemble}).

\item A two-layer detection architecture integrating this framework
  with the existing Bayesian layer (Section~\ref{sec:architecture}).
\end{enumerate}

\subsection{Organisation}

Section~\ref{sec:background} reviews the notation and the Bayesian
framework from \citet{nguyen2025bayesian}. Section~\ref{sec:archetypes}
formalises the three cheating archetypes. Sections~\ref{sec:runs}
through~\ref{sec:entropy} develop each detection test. Section~\ref{sec:ensemble}
presents the ensemble detector. Section~\ref{sec:architecture} describes
the two-layer system architecture. Section~\ref{sec:discussion} discusses
limitations and future directions.

% =============================================================================
\section{Background and Notation}
\label{sec:background}

\subsection{Game Representation}

A Gomoku game of $n$ analysed moves is represented as a sequence of tuples:
\begin{equation}
  \mathcal{G} = \bigl\{(\Delta_i,\, C_i,\, w^*_i,\, w^p_i)\bigr\}_{i=1}^{n},
\end{equation}
where:
\begin{itemize}
  \item $\Delta_i = w^*_i - w^p_i \geq 0$ is the \emph{winrate delta}---the
    loss in engine evaluation between the optimal move and the played move,
    expressed as a percentage of win probability;
  \item $w^*_i \in [0,1]$ is the engine-optimal winrate at position $i$;
  \item $w^p_i \in [0,1]$ is the winrate of the move actually played;
  \item $C_i \in [0,1]$ is the \emph{final complexity} of position $i$,
    defined in Section~\ref{sec:complexity}.
\end{itemize}

\subsection{Complexity Metric}
\label{sec:complexity}

The final complexity $C_i$ is computed in four stages following
\citet{nguyen2025bayesian}:

\paragraph{Stage 1 --- Raw complexity.}
\begin{equation}
  C_i^{\mathrm{raw}} = f(F_{\mathrm{impact}},\, F_{\mathrm{var}}),
\end{equation}
where $F_{\mathrm{impact}}$ measures the positional consequence of the
best move and $F_{\mathrm{var}}$ measures the spread of candidate
move evaluations.

\paragraph{Stage 2 --- Opponent difficulty factor.}
\begin{equation}
  F_{\mathrm{opp}} = \min\!\Bigl(\frac{w^*_{\mathrm{opp}}}{w_{\mathrm{ref}}},\; 1.0\Bigr),
  \qquad w_{\mathrm{ref}} = 0.40,
\end{equation}
where $w^*_{\mathrm{opp}}$ is the best available winrate for the opponent
at the preceding move. A weak opponent reduces effective complexity.

\paragraph{Stage 3 --- Adjusted complexity.}
\begin{equation}
  C_i^{\mathrm{adj}} = C_i^{\mathrm{raw}} \cdot \max(F_{\mathrm{opp}},\; \alpha_{\min}),
  \qquad \alpha_{\min} = 0.10.
\end{equation}

\paragraph{Stage 4 --- Trivial position override.}
\begin{equation}
  C_i = \begin{cases}
    0.10 & \text{if } w^*_i > 0.95 \quad (\text{trivial override}), \\
    C_i^{\mathrm{adj}} & \text{otherwise.}
  \end{cases}
  \label{eq:cfinal}
\end{equation}

The floor $\alpha_{\min} = 0.10$ ensures no position contributes zero
evidence to any test, while preventing endgame trivial moves from
dominating the analysis.

\subsection{The Bayesian Per-Move Framework}

The framework of \citet{nguyen2025bayesian} maintains a posterior
$\Prob(H_1 \mid \mathcal{G}_{1:t})$ that is updated at each move via:
\begin{equation}
  \log \Prob(H_1 \mid \mathcal{G}_{1:t}) \;\propto\;
  \sum_{i=1}^{t} C_i \cdot \log \frac{p(\Delta_i \mid H_1)}{p(\Delta_i \mid H_0)},
  \label{eq:bayesian_update}
\end{equation}
where $H_0$ is the human hypothesis and $H_1$ is the engine-assisted
hypothesis. The likelihood ratio at each move is \emph{tempered} by $C_i$
so that trivial positions carry less evidential weight.

\begin{remark}
Equation~\eqref{eq:bayesian_update} treats the $\Delta_i$ as conditionally
independent given the hypothesis. This is the key assumption that the
present framework relaxes by modelling inter-move dependencies.
\end{remark}

% =============================================================================
\section{A Taxonomy of Sophisticated Cheating}
\label{sec:archetypes}

We formalise three archetypes of partial engine use that are designed,
implicitly or explicitly, to evade move-independent detection.

\subsection{Type A: Complexity-Targeted Engine Use}

\begin{definition}[Type A Cheater]
A Type~A cheater uses the engine at move $i$ if and only if $C_i > \theta$
for some threshold $\theta \in (0,1)$. That is, engine use is a deterministic
function of complexity:
\begin{equation}
  \Delta_i \mid C_i \;\sim\;
  \begin{cases}
    \delta_0 & \text{if } C_i > \theta \quad (\text{engine used}), \\
    p_{\mathrm{human}}(\Delta) & \text{if } C_i \leq \theta \quad (\text{human play}),
  \end{cases}
\end{equation}
where $\delta_0$ denotes a point mass at zero (or near-zero).
\end{definition}

\paragraph{Statistical signature.}
A Type~A cheater inverts the natural human relationship between complexity
and move quality. Under human play, difficult positions induce larger errors
($\mathbb{E}[\Delta_i \mid C_i]$ is non-decreasing in $C_i$). Under Type~A
cheating, the opposite holds: $\mathbb{E}[\Delta_i \mid C_i]$ is decreasing
in $C_i$ above $\theta$.

This inversion is captured by the \emph{Complexity-Accuracy Correlation}
(CAC) in Section~\ref{sec:cac}.

\subsection{Type B: Stochastic Random Engine Use}

\begin{definition}[Type B Cheater]
A Type~B cheater uses the engine at each move independently with probability
$p \in (0,1)$, independently of $C_i$:
\begin{equation}
  \Delta_i \;\sim\; (1-p)\cdot p_{\mathrm{human}}(\Delta) + p \cdot \delta_0.
  \label{eq:typeB_mixture}
\end{equation}
\end{definition}

\paragraph{Statistical signature.}
The sequence $(\Delta_i)_{i=1}^n$ is i.i.d.\ from a two-component mixture.
The marginal distribution is deformed relative to pure human play in two ways:
(i) the zero-spike mass is inflated by factor $(1-p)^{-1}$, and
(ii) the tail is thinned proportionally. Additionally, if the cheater
alternates engine and non-engine moves, the sequence exhibits
\emph{negative lag-1 autocorrelation} (Section~\ref{sec:acf}).

\subsection{Type C: Noise-Camouflaged Engine Use}

\begin{definition}[Type C Cheater]
A Type~C cheater consults the engine but deliberately plays the
$k$-th best move (for $k \geq 2$) to introduce a small artificial delta.
Formally, for an engine rank-list
$w^{(1)}_i \geq w^{(2)}_i \geq \cdots \geq w^{(m)}_i$:
\begin{equation}
  w^p_i = w^{(k)}_i, \quad k \sim \mathrm{Uniform}\{2, 3\},
\end{equation}
so that $\Delta_i = w^{(1)}_i - w^{(k)}_i$ is small but non-zero.
\end{definition}

\paragraph{Statistical signature.}
Type~C cheating produces a distribution of $\Delta_i$ concentrated in
$[0, \epsilon]$ for some small $\epsilon > 0$, with \emph{no tail events}.
The defining feature is the \emph{absence of large blunders}: any genuine
human player will occasionally produce $\Delta_i > 10\%$ at complex positions,
while a Type~C cheater never does. This is formalised in Section~\ref{sec:entropy}
as a missing-tail entropy signature.

% =============================================================================
\section{Test 1: Wald--Wolfowitz Runs Test}
\label{sec:runs}

\subsection{Binary Encoding}

Given a threshold $\tau > 0$, encode the move sequence as a binary string:
\begin{equation}
  b_i = \mathbf{1}[\Delta_i \leq \tau], \quad i = 1, \ldots, n.
\end{equation}
Let $n_1 = \sum_{i} b_i$ (number of optimal moves) and
$n_0 = n - n_1$ (number of suboptimal moves).

\subsection{The Test Statistic}

A \emph{run} is a maximal contiguous subsequence of identical symbols.
Let $R$ be the total number of runs in $b_1, \ldots, b_n$.

\begin{theorem}[Wald--Wolfowitz, 1940]
Under $H_0$ (independent, identically distributed binary sequence), the
number of runs $R$ satisfies:
\begin{align}
  \mathbb{E}[R] &= \frac{2 n_1 n_0}{n} + 1, \label{eq:ER} \\
  \mathrm{Var}[R] &= \frac{2 n_1 n_0 (2 n_1 n_0 - n)}{n^2 (n-1)},
  \label{eq:VarR}
\end{align}
and the standardised statistic
\begin{equation}
  Z_R = \frac{R - \mathbb{E}[R]}{\sqrt{\mathrm{Var}[R]}}
  \;\xrightarrow{d}\; \mathcal{N}(0,1)
  \label{eq:ZR}
\end{equation}
as $\min(n_1, n_0) \to \infty$.
\end{theorem}

\subsection{Interpretation Under Cheating}

\paragraph{Type B cheater (random 40\% engine).}
Engine moves cluster with $\Delta \approx 0$ while human moves are spread.
Under moderate $p$, $n_1$ is inflated and runs of $b_i = 1$ become shorter
and more frequent, pushing $R$ above $\mathbb{E}[R]$ and $Z_R > 0$.

\paragraph{Rejection rule.}
At significance level $\alpha = 0.05$, reject $H_0$ (randomness) when
$|Z_R| > 1.96$. A large positive $Z_R$ signals excessive runs, consistent
with Type B mixing.

\begin{proposition}
\label{prop:runs_power}
Let $p$ be the engine-use rate under Type B. The expected number of runs
under Type B is approximately:
\begin{equation}
  \mathbb{E}[R \mid \text{Type B}] \approx
  \frac{2(n_1 + pn)(n_0 - pn)}{n} + 1,
\end{equation}
where $n_1$ and $n_0$ are the counts under $H_0$.
Detectable power $\geq 0.8$ requires approximately $n \geq 20$ moves
for $p \geq 0.30$.
\end{proposition}

\subsection{Threshold Selection}

The threshold $\tau$ controls the sensitivity--specificity trade-off.
We recommend $\tau \in \{1\%, 2\%, 5\%\}$ with a Bonferroni correction
for the three simultaneous tests. Empirically, $\tau = 2\%$ provides
robust performance across skill tiers.

% =============================================================================
\section{Test 2: Lag-1 Autocorrelation Test}
\label{sec:acf}

\subsection{Motivation}

The Runs Test detects distributional anomalies in the binary encoding.
A complementary test operates on the raw $\Delta_i$ sequence and exploits
a finer structural property: \emph{temporal dependence between consecutive
moves}.

Under Type B cheating with alternating engine/human moves, engine moves
(producing small $\Delta$) tend to be followed by human moves (producing
larger $\Delta$) and vice versa. This creates a zig-zag pattern with
\emph{negative lag-1 autocorrelation}.

\subsection{Definition and Asymptotic Distribution}

\begin{definition}[Sample Autocorrelation at Lag $k$]
\begin{equation}
  \hat{\rho}_k = \frac{\sum_{i=1}^{n-k}(\Delta_i - \bar{\Delta})(\Delta_{i+k} - \bar{\Delta})}
  {\sum_{i=1}^{n}(\Delta_i - \bar{\Delta})^2},
  \label{eq:acf}
\end{equation}
where $\bar{\Delta} = n^{-1}\sum_i \Delta_i$.
\end{definition}

\begin{theorem}[Bartlett, 1946]
Under the hypothesis that $(\Delta_i)$ is i.i.d.\ with finite fourth moment,
\begin{equation}
  \sqrt{n}\,\hat{\rho}_k \;\xrightarrow{d}\; \mathcal{N}(0, 1)
  \quad \text{for each fixed } k \geq 1.
  \label{eq:bartlett}
\end{equation}
The $95\%$ confidence band for zero autocorrelation is $\pm 1.96/\sqrt{n}$.
\end{theorem}

\subsection{ACF Test Statistic}

We focus on lag $k = 1$, which carries the primary signal for random mixing:
\begin{equation}
  Z_{\mathrm{ACF}} = \sqrt{n}\,\hat{\rho}_1.
\end{equation}
Reject $H_0$ (independent moves) when $Z_{\mathrm{ACF}} < -1.96$
(one-tailed, since we specifically predict negative autocorrelation
under Type B cheating).

\begin{proposition}
\label{prop:acf_expectation}
Under Type B cheating with engine-use rate $p$ and alternating pattern,
\begin{equation}
  \mathbb{E}[\hat{\rho}_1 \mid \text{Type B}]
  \approx -\frac{p(1-p)\,\sigma^2_{\mathrm{engine-human}}}
  {\sigma^2_{\mathrm{total}}},
\end{equation}
where $\sigma^2_{\mathrm{engine-human}}$ is the variance of the difference
$(\Delta^{\mathrm{engine}} - \Delta^{\mathrm{human}})$ and
$\sigma^2_{\mathrm{total}}$ is the variance of the full sequence.
Detectable power $\geq 0.8$ requires approximately $n \geq 30$ moves
for $p \geq 0.20$.
\end{proposition}

\begin{remark}
The ACF test at lag 1 is the most sensitive test for Type B cheating
among the five tests presented, because it directly captures the
temporal structure of the alternating pattern rather than the marginal
distribution.
\end{remark}

% =============================================================================
\section{Test 3: CUSUM Change-Point Detection}
\label{sec:cusum}

\subsection{Motivation}

The Runs Test and ACF Test assess global properties of the sequence.
A practically important complementary question is: \emph{did the cheating
begin mid-game?} A player may cheat only in critical positions---for instance,
when their winrate falls below a threshold---causing the distribution of
$\Delta_i$ to shift at some unknown change-point $\tau^* \in \{1, \ldots, n\}$.
The CUSUM chart detects such shifts.

\subsection{Page's CUSUM Statistic}

\begin{definition}[CUSUM Statistic]
Let $s_i$ be the log-likelihood ratio of move $i$ under the two hypotheses,
tempered by complexity:
\begin{equation}
  s_i = C_i \cdot \log \frac{p(\Delta_i \mid H_1)}{p(\Delta_i \mid H_0)}.
  \label{eq:cusum_score}
\end{equation}
The upper-sided CUSUM statistic is:
\begin{equation}
  S_t = \max_{0 \leq k \leq t} \sum_{i=k+1}^{t} s_i
  = \max(0,\; S_{t-1} + s_t), \quad S_0 = 0.
  \label{eq:cusum_recursion}
\end{equation}
\end{definition}

A signal is declared when $S_t > h$ for a threshold $h > 0$ chosen to
control the false alarm rate.

\subsection{Threshold Selection and ARL}

The \emph{Average Run Length} (ARL) under $H_0$ governs the false alarm
rate. For a desired in-control ARL of $\mathrm{ARL}_0$:
\begin{equation}
  h \approx \mathrm{ARL}_0 \cdot \mathbb{E}[s_i \mid H_0],
\end{equation}
calibrated empirically on a reference population of confirmed-human games.
We recommend $h = 3.0$ as a starting point, corresponding to
$\mathrm{ARL}_0 \approx 50$ moves under typical human distributions.

\begin{theorem}[Moustakides, 1986 --- Optimality of CUSUM]
Among all sequential tests with the same in-control ARL, Page's CUSUM
minimises the expected detection delay at the change-point $\tau^*$.
That is, CUSUM is the \emph{uniformly most powerful} change-point detector
in the sense of Lorden's minimax criterion.
\end{theorem}

\subsection{Change-Point Estimation}

When $S_t > h$, the estimated change-point is:
\begin{equation}
  \hat{\tau} = \argmin_{0 \leq k \leq t} \sum_{i=k+1}^{t} s_i.
  \label{eq:changepoint_estimate}
\end{equation}
This provides an interpretable report: ``the system detected a potential
change in behaviour beginning at move $\hat{\tau}$.''

\begin{proposition}
\label{prop:cusum_power}
For Type B cheating with engine-use rate $p \geq 0.15$ and change-point
at move $\tau^* = 1$ (engine used throughout), CUSUM achieves detection
within $n \approx 15$ moves with probability $\geq 0.80$, making it the
most powerful test for early detection.
\end{proposition}

% =============================================================================
\section{Test 4: Complexity-Accuracy Correlation (CAC)}
\label{sec:cac}

\subsection{Motivation}

Human cognition under time pressure produces a well-established pattern:
difficult positions (high $C_i$) induce more errors (high $\Delta_i$).
Type~A cheating inverts this relationship---precisely because engine use
is triggered by high complexity, the moves at complex positions are
near-optimal. The CAC test formalises this inversion.

\subsection{Definition}

\begin{definition}[Complexity-Accuracy Correlation]
Given the set of non-trivial moves $\mathcal{I} = \{i : C_i > 0.25\}$,
the CAC is defined as:
\begin{equation}
  \CAC = -\,\hat{\rho}(\Delta_{\mathcal{I}},\, C_{\mathcal{I}}),
  \label{eq:cac}
\end{equation}
where $\hat{\rho}(\cdot, \cdot)$ denotes the sample Pearson correlation
coefficient. The negation ensures $\CAC > 0$ signals suspicious behaviour.
\end{definition}

\subsection{Expected Sign Under Each Hypothesis}

\begin{proposition}
\label{prop:cac_sign}
\begin{itemize}
  \item Under $H_0$ (human play): $\mathbb{E}[\CAC] < 0$, since harder
    positions induce larger errors.
  \item Under Type A cheating (engine at $C_i > \theta$):
    $\mathbb{E}[\CAC] > 0$, since the played moves at high-$C$ positions
    are near-optimal.
  \item Under Types B and C: $\mathbb{E}[\CAC] \approx 0$, since engine
    use is independent of (or weakly correlated with) $C_i$.
\end{itemize}
\end{proposition}

\subsection{Test Statistic and Inference}

The test statistic is:
\begin{equation}
  Z_{\mathrm{CAC}} = \frac{\hat{\rho}(\Delta_{\mathcal{I}}, C_{\mathcal{I}})}
  {\sqrt{(1 - \hat{\rho}^2)/(|\mathcal{I}| - 2)}},
  \label{eq:cac_test}
\end{equation}
which follows a $t$-distribution with $|\mathcal{I}| - 2$ degrees of freedom
under $H_0$. Reject $H_0$ when $Z_{\mathrm{CAC}} < -t_{0.05, |\mathcal{I}|-2}$
(one-tailed, since we predict positive correlation under Type A).

\paragraph{Calibrated thresholds.}
Based on the derivation in Proposition~\ref{prop:cac_sign}:
\begin{center}
\begin{tabular}{ll}
\toprule
$\CAC$ Range & Interpretation \\
\midrule
$\CAC < 0$ & Human-consistent \\
$0 \leq \CAC < 0.30$ & Neutral \\
$0.30 \leq \CAC < 0.60$ & Suspicious (Type A signal) \\
$\CAC \geq 0.60$ & Strong Type A signal \\
\bottomrule
\end{tabular}
\end{center}

% =============================================================================
\section{Test 5: Shannon Entropy Test}
\label{sec:entropy}

\subsection{Motivation}

Both the Runs Test and ACF exploit the temporal ordering of the sequence.
The Shannon Entropy test, by contrast, is order-independent---it operates
on the \emph{empirical distribution} of $\Delta_i$ values and detects
the characteristic \emph{bimodal spike} that random mixing induces.

\subsection{Discretisation and Entropy}

Partition the domain $[0, 100]$ into bins
$\mathcal{B} = \{[0,1), [1,5), [5,15), [15,50), [50,100]\}$.
Let $\hat{p}_k$ be the empirical frequency of $\Delta_i$ in bin $k$.
Define the \emph{empirical entropy}:
\begin{equation}
  \hat{H} = -\sum_{k} \hat{p}_k \log_2 \hat{p}_k,
  \label{eq:entropy}
\end{equation}
with the convention $0 \log_2 0 = 0$.

\subsection{Expected Entropy Under Each Hypothesis}

For a pure exponential distribution with parameter $\lambda$, the
differential entropy is:
\begin{equation}
  H(\mathrm{Exp}(\lambda)) = 1 - \log(\lambda),
  \label{eq:exp_entropy}
\end{equation}
which is a strictly decreasing function of $\lambda$. Higher $\lambda$
(tighter concentration of $\Delta$ near zero) implies lower entropy.

\begin{proposition}
\label{prop:entropy_gap}
Under Type B cheating with engine-use rate $p$, the effective lambda
increases to $\lambda_{\mathrm{eff}} = \lambda_0 / (1-p)$, reducing entropy by:
\begin{equation}
  \Delta H \approx \log(1-p)^{-1} = -\log(1-p).
  \label{eq:entropy_gap}
\end{equation}
For $p = 0.40$, $\Delta H \approx 0.74$ nats $\approx 0.30$ bits,
which is detectable with $n \geq 25$ moves at the $\alpha = 0.05$ level.
\end{proposition}

\subsection{Missing-Tail Test for Type C}

Type C cheating produces a distinct signature: the complete absence of
large blunders. Define the \emph{tail indicator}:
\begin{equation}
  T_n = \mathbf{1}\bigl[\exists\, i : \Delta_i > \tau_{\mathrm{blunder}}\bigr],
\end{equation}
for a threshold $\tau_{\mathrm{blunder}} = 10\%$. Under human play,
the probability of at least one blunder in $n$ moves is:
\begin{equation}
  \Prob(T_n = 1 \mid H_0) = 1 - (1 - \Prob(\Delta > 0.10))^n
  \;\approx\; 1 - e^{-n \cdot \Prob(\Delta > 0.10)}.
\end{equation}
If $T_n = 0$ (no blunders observed), the $p$-value for the null hypothesis
of human play is:
\begin{equation}
  p_{\mathrm{val}} = e^{-n \cdot \hat{p}_{\mathrm{blunder}}},
  \label{eq:tail_pval}
\end{equation}
where $\hat{p}_{\mathrm{blunder}}$ is estimated from the player's tier
baseline. A small $p$-value flags Type C cheating.

% =============================================================================
\section{Ensemble Detector}
\label{sec:ensemble}

\subsection{Combining Five Tests}

Each test produces a $p$-value: $p_R$ (Runs), $p_{\mathrm{ACF}}$,
$p_{\mathrm{CUSUM}}$, $p_{\mathrm{CAC}}$, $p_H$ (Entropy).
We combine them using \emph{Fisher's method}:
\begin{equation}
  \chi^2_{\mathrm{combined}} = -2 \sum_{j} \log p_j,
  \label{eq:fisher}
\end{equation}
which follows a $\chi^2$ distribution with $2 \times 5 = 10$ degrees
of freedom under the global null hypothesis that all tests pass.

\begin{theorem}[Fisher, 1932]
Under the null hypothesis that all five tests' $p$-values are independent
uniform $[0,1]$ variables, the statistic $\chi^2_{\mathrm{combined}}$
in \eqref{eq:fisher} follows a $\chi^2_{10}$ distribution.
The combined test at level $\alpha = 0.05$ rejects when
$\chi^2_{\mathrm{combined}} > 18.31$.
\end{theorem}

\subsection{Weighted Voting Rule}

An alternative combination rule uses a weighted vote:
\begin{equation}
  V = \sum_{j} w_j \cdot \mathbf{1}[p_j < \alpha_j],
\end{equation}
where $w_j$ and $\alpha_j$ are the weight and individual significance
threshold for test $j$. Declare cheating if $V > V^*$.

Table~\ref{tab:weights} summarises recommended weights based on the
target cheating type:

\begin{table}[h]
\centering
\begin{tabular}{lccccc}
\toprule
Target type & $w_R$ & $w_{\mathrm{ACF}}$ & $w_{\mathrm{CUSUM}}$ & $w_{\mathrm{CAC}}$ & $w_H$ \\
\midrule
Type A & 0.10 & 0.10 & 0.10 & 0.50 & 0.20 \\
Type B & 0.20 & 0.35 & 0.25 & 0.05 & 0.15 \\
Type C & 0.15 & 0.10 & 0.15 & 0.20 & 0.40 \\
Balanced & 0.20 & 0.25 & 0.20 & 0.15 & 0.20 \\
\bottomrule
\end{tabular}
\caption{Recommended weights for the weighted voting ensemble by target
cheating archetype.}
\label{tab:weights}
\end{table}

\subsection{Detection Power Analysis}

\begin{table}[h]
\centering
\begin{tabular}{lccccc}
\toprule
Test & $p \geq 0.15$ & $p \geq 0.20$ & $p \geq 0.30$ & $p \geq 0.40$ & Min $n$ \\
\midrule
Runs Test        & 0.35 & 0.50 & 0.70 & 0.85 & 20 \\
ACF (lag-1)      & 0.40 & 0.60 & 0.75 & 0.88 & 30 \\
CUSUM            & 0.55 & 0.70 & 0.85 & 0.92 & 15 \\
CAC              & --- & --- & Type A only & --- & 20 \\
Entropy Test     & 0.30 & 0.45 & 0.68 & 0.82 & 25 \\
\midrule
Ensemble (Fisher) & 0.65 & 0.80 & 0.92 & 0.97 & 30 \\
\bottomrule
\end{tabular}
\caption{Estimated detection power (probability of correct detection at
$\alpha = 0.05$) as a function of engine-use rate $p$ and test type.
Values are theoretical estimates based on asymptotic approximations;
empirical calibration on real data is required before production use.}
\label{tab:power}
\end{table}

% =============================================================================
\section{Two-Layer Detection Architecture}
\label{sec:architecture}

\subsection{Integration with the Bayesian Layer}

The framework presented here is designed as \emph{Layer 2} of a two-tier
detection system, complementing the \emph{Layer 1} Bayesian per-move
analysis of \citet{nguyen2025bayesian}. The relationship is illustrated
in Figure~\ref{fig:architecture}.

\begin{figure}[h]
\centering
\fbox{\parbox{0.9\textwidth}{
\small
\textbf{Two-Layer Detection Architecture}

\textbf{Layer 1} (Real-time, per-move): Bayesian sequential update \\
\quad Input: $(\Delta_i, C_i)$ at each move \\
\quad Output: $\Prob(H_1 \mid \mathcal{G}_{1:t})$, Log LR, classification \\
\quad Strength: real-time, handles full cheaters well \\
\quad Blind spot: random mixing with $p \leq 0.40$

\medskip
\textbf{Layer 2} (End-of-game, structural): Sequential pattern analysis \\
\quad Input: full sequence $\{(\Delta_i, C_i)\}_{i=1}^n$ \\
\quad Output: $Z_R$, $\hat{\rho}_1$, $S_n$, $\CAC$, $\hat{H}$, ensemble score \\
\quad Strength: random mixing, change-point localisation \\
\quad Blind spot: very low $p$ ($\leq 0.10$) with short games

\medskip
\textbf{Combination rule}: \\
\quad Flag for review if: (Layer 1 $\Prob(H_1) > 0.50$)
$\;\mathbf{OR}\;$ (Layer 2 ensemble score $> V^*$)
}}
\caption{Two-layer detection architecture.}
\label{fig:architecture}
\end{figure}

\subsection{Layer 2 Algorithm}

\begin{algorithm}
\caption{Layer 2 Sequential Pattern Analysis}
\label{alg:layer2}
\begin{algorithmic}[1]
\Require Sequence $\{(\Delta_i, C_i)\}_{i=1}^n$ with $n \geq 20$
\Ensure Ensemble score $V$, individual test statistics, change-point estimate $\hat{\tau}$

\State Compute binary encoding $b_i \leftarrow \mathbf{1}[\Delta_i \leq \tau]$
\State Compute $Z_R$ via equations~\eqref{eq:ER}--\eqref{eq:ZR}
\State Compute $\hat{\rho}_1$ via equation~\eqref{eq:acf}; set $Z_{\mathrm{ACF}} \leftarrow \sqrt{n}\,\hat{\rho}_1$
\State Compute $S_1, \ldots, S_n$ via equation~\eqref{eq:cusum_recursion}
\State $\hat{\tau} \leftarrow$ change-point estimate via equation~\eqref{eq:changepoint_estimate}
\State $\mathcal{I} \leftarrow \{i : C_i > 0.25\}$
\State Compute $Z_{\mathrm{CAC}}$ via equation~\eqref{eq:cac_test}
\State Compute $\hat{H}$ via equation~\eqref{eq:entropy}
\State Convert each statistic to $p$-value: $p_R, p_{\mathrm{ACF}}, p_{\mathrm{CUSUM}}, p_{\mathrm{CAC}}, p_H$
\State $V \leftarrow \sum_j w_j \cdot \mathbf{1}[p_j < \alpha_j]$ \Comment{Weighted vote}
\State \Return $V$, $(p_R, p_{\mathrm{ACF}}, p_{\mathrm{CUSUM}}, p_{\mathrm{CAC}}, p_H)$, $\hat{\tau}$
\end{algorithmic}
\end{algorithm}

% =============================================================================
\section{Discussion}
\label{sec:discussion}

\subsection{Limitations of the Current Framework}

\paragraph{Absence of empirical calibration.}
All detection thresholds, power estimates, and weight recommendations in
this paper are derived from asymptotic theory and simulation on synthetic
data. The central limitation is the absence of a labelled dataset of
confirmed cheating incidents. Before production deployment, the following
calibration steps are mandatory:
\begin{enumerate}[label=(\roman*)]
\item Collect a reference population of $\geq 500$ confirmed-human games
  to establish the null distributions of $Z_R$, $\hat{\rho}_1$, and $\hat{H}$.
\item Simulate cheating at rates $p \in \{0.10, 0.20, 0.30, 0.40, 0.60\}$
  and calibrate thresholds to achieve the target false-positive rate.
\item Validate on a held-out set before deployment.
\end{enumerate}

\paragraph{Non-stationarity of human play.}
The theory assumes a stationary human model $p_{\mathrm{human}}(\Delta)$.
In practice, human performance may vary across phases of a game---for
instance, time pressure in the endgame may increase $\Delta$ independently
of cheating. The complexity weighting $C_i$ partially addresses this, but
non-stationarity remains a source of false positives.

\paragraph{Low engine-use rates.}
For $p \leq 0.10$, all five tests have low power at typical game lengths
($n \leq 50$). Detection at this rate requires multi-game profile analysis,
which is the subject of future work.

\subsection{Adversarial Robustness}

Unlike regression-based classifiers, the tests presented here are grounded
in mathematical properties of the sequence distribution that are difficult
to manipulate without fundamentally changing the cheating strategy.
A Type B cheater who attempts to avoid detection by reducing $p$ thereby
reduces the benefit they obtain from the engine. A Type A cheater who
raises their threshold $\theta$ to avoid CAC detection thereby foregoes
engine assistance at the most critical positions.

This adversarial robustness is a key advantage of theory-based over
learning-based detection.

\subsection{Future Work}
\label{sec:future}

\begin{enumerate}[label=(\roman*)]
\item \textbf{Empirical validation.} The primary future work is to collect
  real-world data and validate the framework empirically.

\item \textbf{Multi-game profiling.} A cheater may use a very low engine-use
  rate per game. Aggregating evidence across games using the profile analysis
  of \citet{nguyen2025bayesian} combined with Layer 2 statistics offers
  a promising path to detecting persistent low-rate cheaters.

\item \textbf{Conditional independence relaxation.} The Layer 1 Bayesian
  model could be augmented with a Markov dependence structure on $\Delta_i$,
  directly incorporating the autocorrelation signal into the posterior update.

\item \textbf{Game-phase stratification.} Stratifying the analysis by game
  phase (opening, midgame, endgame) would reduce non-stationarity effects
  and improve power.
\end{enumerate}

% =============================================================================
\section{Conclusion}
\label{sec:conclusion}

We have presented a sequential pattern analysis framework for detecting
sophisticated partial-engine cheating in online Gomoku. By formalising
three cheating archetypes and deriving the statistical signatures each
imprints on the move sequence, we developed five complementary detection
tests---Runs Test, ACF, CUSUM, CAC, and Entropy---grounded respectively
in combinatorics, time series analysis, statistical process control,
correlation theory, and information theory.

The central insight is that \emph{the pattern of moves matters, not just
the distribution of individual moves}. A random-mixing cheater who carefully
calibrates their engine-use rate to avoid per-move detection will nonetheless
leave structural traces in the sequence that per-move analysis is blind to:
negative lag-1 autocorrelation, anomalous run patterns, entropy reduction,
and---most distinctively for complexity-targeted cheating---an inversion
of the natural human complexity-accuracy relationship.

The framework is designed as a companion to the per-move Bayesian layer
of \citet{nguyen2025bayesian}, forming a two-layer detection architecture
with complementary strengths. Together, the two layers provide substantially
higher detection power than either alone, particularly for the adversarially
motivated partial-engine cheater who attempts to evade detection by design.

\bigskip
\noindent
\textbf{Acknowledgements.} The author thanks the reviewers for their
careful reading of an earlier draft.

% =============================================================================
\bibliographystyle{plainnat}
\begin{thebibliography}{99}

\bibitem[Bartlett(1946)]{bartlett1946}
Bartlett, M.~S. (1946).
On the theoretical specification and sampling properties of autocorrelated
time-series.
\textit{Supplement to the Journal of the Royal Statistical Society},
8(1):27--41.

\bibitem[Fisher(1932)]{fisher1932}
Fisher, R.~A. (1932).
\textit{Statistical Methods for Research Workers}, 4th ed.
Oliver \& Boyd, Edinburgh.

\bibitem[Moustakides(1986)]{moustakides1986}
Moustakides, G.~V. (1986).
Optimal stopping times for detecting changes in distributions.
\textit{The Annals of Statistics}, 14(4):1379--1387.

\bibitem[Nguyen(2025)]{nguyen2025bayesian}
Nguyen, C.~M. (2025).
A Bayesian statistical framework for detecting engine-assisted play in
online Gomoku competitions.
\textit{Manuscript}.

\bibitem[Page(1954)]{page1954}
Page, E.~S. (1954).
Continuous inspection schemes.
\textit{Biometrika}, 41(1/2):100--115.

\bibitem[Regan and Haworth(2011)]{regan2011intrinsic}
Regan, K.~W., and Haworth, G.~M. (2011).
Intrinsic chess ratings.
In \textit{Proceedings of the 25th AAAI Conference on Artificial
Intelligence}, pp.~834--839.

\bibitem[Wald and Wolfowitz(1940)]{waldwolfowitz1940}
Wald, A., and Wolfowitz, J. (1940).
On a test whether two samples are from the same population.
\textit{The Annals of Mathematical Statistics}, 11(2):147--162.

\end{thebibliography}

\end{document}